\section{The surplus construction and its execution certificate}
\label{imp:section}

This section specifies the external algorithmic input and derives the
history properties used below.  The input is the construction of
Cohen-Addad, Grandoni, Lee, Schwiegelshohn, and Svensson
\cite{CGLSS2026}, in the version arXiv:2503.10972v2 dated 19 May 2026.
We use its logarithmic-surplus construction and its robust dual analysis.
The strictness proof below will be applied to the resulting certificate.

Throughout this section, clients and facilities are distinct indexed points,
all distances between distinct original indices are positive integers bounded by a fixed
polynomial in $N=n+|F|$, and $n>0$.  Fix a sufficiently small rational
$v\in(0,1/6)$ and put
\[
 \delta=3v,\qquad K=1+v^2,\qquad \lambda=2f>0,
 \qquad [x]^+=\max\{x,0\}.
\]
The parameter $v$ is the small parameter denoted by $\varepsilon$ in the
source's Section~3; the final approximation slack in its Theorem~1.2 need
not have the same value.

\subsection{The sequence data and the external result}

A free copy $\widetilde i$ of a regular facility $i\in F$ has a specified
length $u(\widetilde i)\ge0$.  Attach it as a leaf at $i$.  Thus
\[
 d(p,\widetilde i)=d(p,i)+u(\widetilde i),\qquad
 d(\widetilde i,\widetilde j)
   =u(\widetilde i)+d(i,j)+u(\widetilde j)
\]
for distinct copies.  Zero-length leaves may be identified metrically
while retaining their labels for accounting.  In particular all triangle
inequalities used below hold.  Each copy's length is fixed throughout the
execution considered here.

Here is the relevant sequence definition, stated with our notation.
At phase value $z$ and a set $W$ of facilities define
\[
 A_W=\{j\in D:z<(1-\delta)d(j,W)\},\qquad I_W=D\setminus A_W.
\]
The convention $d(j,\varnothing)=+\infty$ is used only when every client
is active.  A regular facility $i$ is $\eta$-openable relative to $(z,W)$
if there are bids $(\tau_j)_{j\in A_W}$ satisfying
\begin{align}
 \tau_j&=z &&(j\in A_W,\ d(i,j)>vz),\label{imp:near-bids}\\
 z\le\tau_j&\le\min\{Kz,(1-\delta)d(j,W)\}
       &&(j\in A_W,\ d(i,j)\le vz),\label{imp:bid-cap}
\end{align}
and the payment and dual conditions
\begin{equation}\label{imp:opening-payment}
 \begin{split}
 \sum_{j\in A_W}[\tau_j-(1-\delta)d(i,j)]^+
  &+\sum_{j\in I_W}[(1-\delta)d(j,W)-(1-\delta)d(i,j)]^+\\
  &\ge\lambda-n\eta,
 \end{split}
\end{equation}
\begin{equation}\label{imp:opening-dual}
 \sum_{j\in A_W}[\tau_j-d(i_0,j)]^+
  +\sum_{j\in I_W}[\tau_a-2d(i_0,j)-d(i_0,a)]^+
   \le\lambda
 \qquad(i_0\in F,\ a\in A_W).
\end{equation}
Testing $i_0\in F$ also tests free copies: replacing $i_0$ by a free
copy only increases the distances subtracted in
\eqref{imp:opening-dual}.  These are Definition~3 of
\cite{CGLSS2026}; $0$-openable means openable with the exact payment
requirement.

Suppose the actual set at the start of a phase is $Q$.  A phase sequence
$H=(a_1,\ldots,a_s)$ is $\eta$-valid if every entry is either a free
copy or a regular facility $\eta$-openable against some witness
\[
 W_t\supseteq Q\cup\{a_1,\ldots,a_{t-1}\},
\]
and no unopened regular facility is $0$-openable against the actual set
$Q\cup\{a_1,\ldots,a_s\}$ after the sequence.  This last requirement
is maximality for \emph{exact} openability, even when $\eta>0$.
A solution of valid sequences consists of such sequences at successive
phase values $1,K,K^2,\ldots$ and ends with all clients inactive.
This is Definition~4 of \cite{CGLSS2026}.

\begin{theorem}[Surplus construction]
\label{imp:cglss}
For the metric range above and fixed $v$, the construction underlying
Theorem~1.2 of \cite{CGLSS2026} can be run in polynomial time at any
target $2\le h\le |F|$.  With precision $\eta=2^{-N}$, its output has
one of the following forms.
\begin{enumerate}[label=(\roman*)]
\item A solution of $\eta$-valid sequences with exactly $h$ distinct
regular facilities and at most $L=O_v(\log N)$ free copies, for some
final price $f>0$ and specified final copy lengths.
\item A solution of $0$-valid sequences containing only regular
facilities, at price $f=N^{-2}$, with at most $h$ facilities.
\end{enumerate}
One common integer $L$ can be used for all targets on an instance.
For every regular opening, the algorithm also returns its witness set
and certified bids.
\end{theorem}

The merging procedure is the external ingredient: Algorithm~3
(\textsc{LogAdaptive}), Definitions~3--4, and Section~3.5
(\textsc{MergeSolutions}) of \cite{CGLSS2026}. Its target counts
regular facilities before projection of free copies. Section~\ref{imp:instantiation}
justifies the price range, precision, interpolation endpoints and copy
lengths used here; these choices are part of the stated construction.

Before that verification, we prove what follows from any valid-sequence
output of either form. Fix its final price and copy lengths and replay
only its returned sequences. The preceding interpolation search is not
part of this history. The next lemma derives all budget inequalities
from this replay, including the final dual and its free-opening case.
Its no-overbidding invariant will also justify the disabling length
in Section~\ref{imp:instantiation}.

\subsection{One execution supplies all the inequalities}

\begin{lemma}[Surplus solution with a compatible history]
\label{imp:surplus-history}
Given valid sequences of either output form in Theorem~\ref{imp:cglss},
with their price, fixed copy lengths, opening witnesses and certified bids, a
polynomial-time replay produces final budgets $(\alpha_p)_{p\in D}$
and a real solution $U\subseteq F$ of size at most $h+L$ with the
following properties.

The history replays the \emph{final} sequences using their fixed final
price and fixed final copy lengths.  Each atomic step opens one facility
or performs one simultaneous phase update.  Its actual pre-step and
post-step open sets satisfy
$Q_t^-\subseteq Q_t^+$, and these sets accumulate throughout the
history. We call a client's transition to inactivity its removal from
the active set. Let $t(p)$ be the step at which $p$ becomes inactive and
$z_p$ the phase value immediately \emph{before} that step.  At this
pre-step state all active clients have budget $z_p$, and
\begin{align}
 z_p&\le(1-\delta)d(p,Q_{t(p)}^-),\label{imp:active-threshold}\\
 \alpha_p&\le Kz_p,\label{imp:final-cap}\\
 (1-\delta)d(p,Q_{t(p)}^+)&\le\alpha_p\le Kz_p.
 \label{imp:post-removal}
\end{align}
The distance on the left of \eqref{imp:post-removal} cannot increase
at any later step.  If $I_t$ is the actual inactive set just before
step $t=t(p)$, then, for every regular facility $i\in F$,
\begin{equation}
 \sum_{j\notin I_t}[z_p-d(i,j)]^+
 +\sum_{j\in I_t}[(1-\delta)d(j,Q_t^-)-d(i,j)]^+
 \le\lambda.
 \label{imp:no-overbid}
\end{equation}
For these same final budgets, the dual inequality is
\begin{equation}
 \sum_{p\in D}[\alpha_p-2d(i,p)]^+\le\lambda
 \qquad(i\in F).
 \label{imp:dual}
\end{equation}

In the normal branch the history opens exactly $h$ paid regular
facilities before projecting free copies, and
\begin{equation}
 (1-\delta)\operatorname{cost}(U)+h(\lambda-N\eta)
       \le\sum_{p\in D}\alpha_p,
 \qquad \eta=2^{-N}.
 \label{imp:payment}
\end{equation}
In the small-price branch, $|U|\le h$, $\lambda=2/N^2$, $\eta=0$,
and, writing $r=|U|$, the stronger exact-payment inequality is
\begin{equation}
 (1-\delta)\operatorname{cost}(U)+r\lambda
       \le\sum_{p\in D}\alpha_p.
 \label{imp:small-payment}
\end{equation}
In particular its paid term can be dropped.  For every nonempty real
comparison solution $H\subseteq F$ one also has
\begin{equation}
 \sum_{p\in D}\alpha_p
       \le |H|\lambda+2\operatorname{cost}(H).
 \label{imp:comparison}
\end{equation}
\end{lemma}

\begin{proof}
Fix the supplied valid sequences, price and copy lengths. This proof
uses their validity. No facility deletion, change of price, or movement
of a free copy during the construction's search is used as a step of
this history. Write $\beta_j$ for
a transient budget, reserving $\alpha_j$ for its final value.
Initialize $Q=\varnothing$, $z=1$, and $\beta_j=1$ for all clients.

\paragraph{Witness sets and actual updates.}
At a regular opening with witness $W\supseteq Q$, one has
$A_W\subseteq A_Q$ and $I_Q\subseteq I_W$.  Use the certified bid
$\tau_j$ for $j\in A_W$ and the unchanged bid $z$ for
$j\in A_Q\setminus A_W$.  Budgets of already inactive clients stay
fixed.  Every actual active client whose budget increases lies within
$vz$ of the newly opened regular facility.  It therefore immediately
becomes inactive, since
\[
 (1-\delta)d(j,Q\cup\{i\})
     \le(1-\delta)vz<z\le\tau_j.
\]
Consequently every client remaining active has its old budget $z$.
At a free opening no budgets increase and the same conclusion holds.
At a phase update, with old value $z$, each active budget becomes
\[
 \beta_j'=\min\{Kz,(1-\delta)d(j,Q)\};
\]
clients reaching the second value become inactive, and all others
have budget $Kz$, the next phase value.  This proves inductively both
the common active budget and the identity $A=A_Q$ at every atomic
boundary.  It gives the strict form of
\eqref{imp:active-threshold} whenever the right side is finite.

When a regular opening inactivates a client, its final bid is at most
$Kz_p$ and at least its new scaled connection distance. Inactivation
during a free opening has final bid exactly $z_p$. At inactivation
during a phase update, the final bid equals the scaled connection
distance and is at most $Kz_p$. In every case the bid is frozen
after removal from the active set.
This proves \eqref{imp:final-cap}--\eqref{imp:post-removal}, using
the old phase value also for a phase-update removal.
Since $Q$ only grows in the fixed metric extension, later connection
distances cannot increase.  Several removals in one atomic step all
use the same pre-step value and pre-step set; an ordering of tied
clients does not insert fictitious intermediate states.

\paragraph{No overbidding at actual atomic boundaries.}
We include the invariant argument corresponding to Claim~10,
Equation~(33), in Appendix~A.2 of \cite{CGLSS2026}.  At a boundary set
$c_j=(1-\delta)d(j,Q)$ and define
\[
 B_i=\sum_{j\in A}[z-d(i,j)]^+
          +\sum_{j\in I}[c_j-d(i,j)]^+.
\]
Initially $B_i=0$ because all client--facility distances are at
least one.  An opening cannot increase any client's term.
Old inactive clients have smaller or unchanged $c_j$; an active
client that stays active keeps $z$; and an active client that becomes
inactive has new $c_j\le z$.  The latter assertion holds even if its
bid increased, because an increased bid is possible only within
$vz$ of the opening facility.  A free opening has the same
monotonicity without bid increases.  This proves preservation at
opening steps, including openings witnessed by supersets.

Consider a phase update after a maximal phase sequence.  For its old
active set put $b_j(t)=\min\{t,c_j\}$ for $z\le t\le Kz$, and define
\[
 B_i(t)=\sum_{j\in A}[b_j(t)-d(i,j)]^+
             +\sum_{j\in I}[c_j-d(i,j)]^+.
\]
At $t=Kz$ this is precisely the invariant after the simultaneous
update, including newly inactive clients.  If some $B_i(Kz)>\lambda$,
continuity gives a first $t_*\in[z,Kz]$ at which some such facility
$i$ has $B_i(t_*)=\lambda$; choose the first crossing among all
regular facilities.  Then $B_{i_0}(t_*)\le\lambda$ for every
$i_0\in F$.  Such an $i$ is unopened, since a facility in $Q$ has
$B_i(t)=0$ for all these $t$.

Offer $\tau_j=b_j(t_*)$ for $d(i,j)\le vz$ and $\tau_j=z$
otherwise.  These bids satisfy the two update restrictions.
For an outside client, $\delta d(i,j)\ge3v^2z\ge v^2z$, so
\[
 z-(1-\delta)d(i,j)\ge b_j(t_*)-d(i,j).
\]
For a nearby client this comparison holds immediately with
$z$ replaced by $b_j(t_*)$.  Also
$[c_j-(1-\delta)d(i,j)]^+\ge[c_j-d(i,j)]^+$ for an inactive client.
Thus these bids pay at least $B_i(t_*)=\lambda$.
Finally, for every $a\in A$, $j\in I$, and $i_0\in F$,
the triangle inequality and $\tau_a\le c_a$ give
\begin{align*}
 \tau_a-d(i_0,a)-2d(i_0,j)
 &\le (1-\delta)d(j,Q)-d(i_0,j)\\
 &=c_j-d(i_0,j).
\end{align*}
Since $\tau_j\le b_j(t_*)$ for active $j$, the left side of
\eqref{imp:opening-dual} is at most $B_{i_0}(t_*)\le\lambda$.
Hence $i$ is $0$-openable against the actual set $Q$, contradicting
maximality.  This proves the invariant through phase updates and
establishes \eqref{imp:no-overbid} at precisely the required states.

\paragraph{Payment with witness supersets.}
The certificate's payment can be checked directly by a potential
argument; this is the accounting in Lemma~A.1 and Equation~(32)
of \cite{CGLSS2026}.  First compare a regular opening's witness
payments with its actual payments.  For
$j\in A_Q\setminus A_W$ we have
$z\ge(1-\delta)d(j,W)$, whence
\[
 [z-(1-\delta)d(i,j)]^+
 \ge[(1-\delta)d(j,W)-(1-\delta)d(i,j)]^+.
\]
For $j\in I_Q$, replacing $W$ by $Q\subseteq W$ only increases
the available saving.  Clients in $A_W$ keep their certified bids.
Therefore \eqref{imp:opening-payment} remains valid against the
\emph{actual} prefix with the extended actual bids.

Let $r_Q$ count regular facilities opened in the actual prefix.
At every boundary consider
\[
 \mathcal P=\sum_{j\in I}\bigl(\beta_j-(1-\delta)d(j,Q)\bigr)
              -r_Q(\lambda-n\eta).
\]
Initially it is zero.  At a regular opening, already inactive
clients contribute their nonnegative connection savings.
For an old active client a positive contribution
$[\tau_j-(1-\delta)d(i,j)]^+$ makes it inactive, and its new slack
$\tau_j-(1-\delta)d(j,Q\cup\{i\})$ is at least that
contribution.  The total increase before the paid-facility term is
therefore at least $\lambda-n\eta$.
At a free opening, old inactive slacks can only increase and new
inactive slacks are nonnegative.  At a phase update, the new inactive
clients have zero slack, while the old terms and $r_Q$ are unchanged.
Thus $\mathcal P\ge0$ always.  At termination, writing $V$ for the
final virtual facilities and $r$ for its regular count, this gives
\begin{equation}
 (1-\delta)\sum_{j\in D}d(j,V)+r(\lambda-n\eta)
           \le\sum_{j\in D}\alpha_j.
 \label{imp:virtual-payment}
\end{equation}

\paragraph{The final dual on this replay.}
We prove \eqref{imp:dual} for the budgets just constructed. The source
defines this execution in Section~3.4 and immediately
before Appendix~A.1. Its final-dual statement is Lemma~A.2, used to
prove Theorem~3.10. Our argument follows Lemmas~A.2--A.3 and
Equations~(34)--(38) of
\cite[Appendix~A.2]{CGLSS2026}, whose regular-opening argument already
handles witness supersets. We give the rows explicitly, including
the case where the first facility that freezes the comparison facility
is free.  That case needs the no-overbidding invariant in place of an
opening-dual condition.

Fix $i\in F$, write $d_j=d(i,j)$, and put
\[
 D^*=\{j\in D:\alpha_j>2d_j\},\qquad m=|D^*|.
\]
There is nothing to prove if $m=0$.  Order $D^*$ by removal step,
breaking ties by nondecreasing final budget $\alpha_j$, and then
arbitrarily between equal budgets.  Write $j\prec k$ and $j\succeq k$
for this order.  This auxiliary order is specific to the present dual
proof; it imposes no restriction on the arbitrary order of simultaneous
removals allowed in the later strictness argument.  In particular it
does not create intermediate states within an atomic step.

An opening of $h$ at phase value $z$ \emph{freezes} $i$ if
$d(i,h)\le(1+3v)z/2$.  We use only the first such opening, if there is
one.  We first establish the following row for each contributor $k$
removed strictly before that opening, or for every $k\in D^*$ if no
opening freezes $i$:
\begin{equation}\label{imp:pre-freeze-row}
 \sum_{\substack{j\in D^*\\j\succeq k}}[\alpha_k-d_j]^+
 +\sum_{\substack{j\in D^*\\j\prec k}}
       [\alpha_k-2d_j-d_k]^+\le\lambda.
\end{equation}
The elementary comparison used below is that, for a nonempty actual
prefix $Q$, $c_j=(1-\delta)d(j,Q)$, and $u\le c_k$,
\begin{equation}\label{imp:triangle-row}
 u-2d_j-d_k\le c_j-d_j.
\end{equation}
Indeed, $d(k,Q)\le d_k+d_j+d(j,Q)$, and multiplication by
$1-\delta\le1$ proves the assertion.  When $Q=\varnothing$ there are
no inactive clients, so none of the applications to inactive clients
below is needed.

Suppose first that $k$ is removed by opening $h$ at phase value $z$,
before $i$ is frozen.  No contributor can have its budget increased by
this opening: such an increase would require $d(h,j)\le vz$, whereas
\[
 2d_j\ge2d(i,h)-2d(h,j)>(1+v)z\ge Kz\ge\alpha_j,
\]
contradicting $j\in D^*$.  Thus $\alpha_k=z$.  Apply no overbidding
at the actual boundary immediately before the opening.  Every
$j\succeq k$ is then active, so its term is exactly
$[\alpha_k-d_j]^+$.  A $j\prec k$ removed in the same opening is
also still active at this boundary, and its term $[z-d_j]^+$ dominates
$[\alpha_k-2d_j-d_k]^+$.  Any remaining $j\prec k$ is already
inactive; actual activity of $k$ gives $\alpha_k=z<c_k$, so
\eqref{imp:triangle-row} gives the required domination for $j$.
Discarding noncontributors proves \eqref{imp:pre-freeze-row}.

Suppose instead that a simultaneous phase update removes $k$, with old
phase value $z$.  Use the no-overbidding invariant at the actual
\emph{post-update} boundary.  The set $Q$ is unchanged and
$\alpha_k=c_k\le Kz$.  Clients still active have budget $Kz$, which
dominates the budget appearing in either term of
\eqref{imp:pre-freeze-row}.  For an inactive $j\prec k$, use
\eqref{imp:triangle-row}.  An inactive $j\succeq k$ must have been
removed in this same update.  Therefore $c_j=\alpha_j\ge\alpha_k$,
where the last inequality is precisely the reason for the
final-budget tie rule.  Its invariant term dominates
$[\alpha_k-d_j]^+$.  This proves \eqref{imp:pre-freeze-row} also for
phase-update removals.

If no opening freezes $i$, sum \eqref{imp:pre-freeze-row} over all
$k\in D^*$.  Replace each positive part by its argument.  Each budget
has coefficient $m$.  If $j$ has rank $q$, its distance occurs with
coefficient $q+2(m-q)+(q-1)=2m-1$.  Hence
\[
 m\sum_{j\in D^*}\alpha_j-(2m-1)\sum_{j\in D^*}d_j
       \le m\lambda,
\]
which implies $\sum_{j\in D^*}(\alpha_j-2d_j)\le\lambda$.

It remains to consider a first freezing opening $h$ at phase value
$z$.  Let $Q,A,I$ be the actual state immediately before it, and let
$b_j$ be the extended actual bid used at this opening: certified bids
for witness-active clients and $z$ for the other active clients if
$h$ is regular, and $b_j=z$ for all $j\in A$ if $h$ is free.
Every $j\in A\cap D^*$ is removed at this opening and has
$\alpha_j=b_j$.  To see this, suppose that $j$ remains active.
Then $z\le b_j<(1-\delta)d(j,h)$.  Every later actual prefix, and
thus every later witness, contains $h$.  Phase updates and certified
bid updates are consequently capped by $(1-\delta)d(j,h)$; unchanged
active bids also satisfy this bound.  It follows that
$\alpha_j\le(1-\delta)d(j,h)$.  But
\[
 2d_j\ge2d(j,h)-(1+3v)z
    >\bigl(2-(1+3v)(1-\delta)\bigr)d(j,h)
     =(1+9v^2)d(j,h)>\alpha_j,
\]
a contradiction.  This argument applies to free and regular openings
and uses only containment of the actual prefix in future witnesses.

Write $R=I\cap D^*$.  These contributors precede the freezing opening,
so they retain all rows \eqref{imp:pre-freeze-row}.
If $h$ is regular, let $W\supseteq Q$ be its witness and set
\[
 P=A_W\cap D^*,\qquad T=(A\setminus A_W)\cap D^*.
\]
For $j\in P$ the certified bid equals $\alpha_j$, whereas for
$j\in T$ the unchanged actual bid gives $\alpha_j=z$.  Restricting
\eqref{imp:opening-dual} to contributors gives, for every $k\in P$,
\begin{equation}\label{imp:witness-freeze-row}
 \sum_{j\in P}[\alpha_j-d_j]^+
 +\sum_{j\in R\cup T}[\alpha_k-2d_j-d_k]^+\le\lambda.
\end{equation}
For each $k\in T$, apply actual no overbidding immediately before the
opening.  The active budget is $z=\alpha_k$, and $k$ is actually
active, so $\alpha_k<c_k$.  Using \eqref{imp:triangle-row} for
inactive clients gives
\begin{equation}\label{imp:unchanged-freeze-row}
 \sum_{j\in P\cup T}[\alpha_k-d_j]^+
 +\sum_{j\in R}[\alpha_k-2d_j-d_k]^+\le\lambda.
\end{equation}
Thus actual-active but witness-inactive contributors use the actual
invariant, without assuming that the two active sets coincide.

If $h$ is free, set $P=\varnothing$ and $T=A\cap D^*$.
The preceding freezing argument shows that every $k\in T$ is removed
by this free opening with $\alpha_k=z$.  The same actual pre-opening
invariant and \eqref{imp:triangle-row} prove
\eqref{imp:unchanged-freeze-row} for each such $k$.  No witness or
opening-dual condition is needed for the free opening.  This supplies
the free first-freezing case omitted by the invocation of an
opening-dual condition in the cited proof.

In either case put $r=|R|$, $p=|P|$, and $t=|T|$, so $m=r+p+t$.
Sum \eqref{imp:pre-freeze-row} for $k\in R$,
\eqref{imp:witness-freeze-row} for $k\in P$, and
\eqref{imp:unchanged-freeze-row} for $k\in T$, replacing every
positive part by its argument.  Each budget has coefficient $m$:
for $R$ and $T$ it occurs $m$ times in its own row, and for $P$ it
occurs in the first sum of each of the $p$ rows indexed by $P$ and
$m-p$ times in its own second sum.  The distance coefficients are
\[
 \begin{array}{c|c|c}
 \text{client class}&\text{coefficient calculation}&\text{coefficient}\\\hline
 j\in R\text{, rank }q\text{ in }R
    &q+2(r-q)+(q-1)+2p+2t&2m-1\\
 j\in T&r+2p+t+r&2m-t\\
 j\in P&r+t+p+(r+t)&2m-p.
 \end{array}
\]
Consequently
\[
 m\sum_{j\in D^*}\alpha_j
 -(2m-1)\sum_{j\in R}d_j
 -(2m-t)\sum_{j\in T}d_j
 -(2m-p)\sum_{j\in P}d_j\le m\lambda.
\]
Every displayed distance coefficient is at most $2m$; subtracting the
remaining nonnegative distance terms and dividing by $m$ proves
\[
 \sum_{j\in D}[\alpha_j-2d(i,j)]^+
   =\sum_{j\in D^*}(\alpha_j-2d_j)\le\lambda.
\]
This is \eqref{imp:dual} for the same final budgets used in the payment
argument.  It also holds for each free copy by distance domination,
although only $i\in F$ is needed here.

\paragraph{Projection and cardinality.}
Project each free copy of $i$ to $i$ and discard duplicate locations,
giving $U\subseteq F$.  For every client,
$d(j,U)\le d(j,V)$, so projection can only decrease cost.
In the normal branch $r=h$ before projection and $|U|\le h+L$;
substituting into \eqref{imp:virtual-payment} and using $n\le N$
proves \eqref{imp:payment}.  Coincident projected locations do not
reduce the number $h$ of paid regular facilities in that inequality.
In the unmerged branch there are no free copies, $r=|U|\le h$,
and $\eta=0$, giving \eqref{imp:small-payment}.
Finally assign every client to its nearest facility of a nonempty
$H$.  For each $i\in H$, the elementary bound
\[
 \sum_{j\text{ assigned to }i}\alpha_j
 \le2\sum_{j\text{ assigned to }i}d(i,j)
       +\sum_{j\in D}[\alpha_j-2d(i,j)]^+
\]
and \eqref{imp:dual} give \eqref{imp:comparison} upon summation.
All statements therefore use one and the same history and budgets.
\end{proof}

\subsection{Instantiating the construction}
\label{imp:instantiation}

We now verify the choices needed for Theorem~\ref{imp:cglss}.
The merging procedure itself is that of \cite[Section~3.5]{CGLSS2026}.
It uses at most three free copies per phase and stops at exactly the
target number of regular facilities. We check its endpoints and
parameter transfers for the metric range and precision used here.

\paragraph{Price range and precision.}
Initialize at $f=N^{-2}$ and retain the source's
upper price $4nM$, where $M$ bounds all original distances.
At this upper endpoint, the first opening in an unmerged execution
satisfies $nKz\ge\lambda=8nM$, since every active bid is at most
$Kz$. Hence $z>M$, that opening makes all clients inactive, and
the total subsequent switching payment is at most $nM<\lambda$.
The upper-endpoint solution therefore has one facility, as required
to bracket $h\ge2$. At the lower endpoint only cardinality is tested:
if the unmerged run has at most $h$ facilities, return it as the
small-price alternative. Section~\ref{sur:section} proves its cost
bound using strictness.

Perform each binary search to parameter difference at most $\eta/2$,
where $\eta=2^{-N}$. A price difference of at most $\eta/2$ changes
$\lambda=2f$ by at most $\eta\le n\eta$; a copy-length difference
of at most $\eta/2$ changes the total payment by at most $n\eta/2$.
Thus the source's parameter-transfer argument supplies the stated
$\eta$-valid sequences, including when $n=1$. The source maintains
a common $\eta$-valid prefix of earlier phases and transfers a current
$0$-valid phase across one changed parameter
\cite[Section~3.5]{CGLSS2026}. Thus the per-opening validity tolerance
is not compounded from phase to phase.

The transfer also preserves the unrelaxed opening-dual condition.
Here is the scaled comparison needed when a client becomes inactive
under a changed copy length, completing the corresponding calculation
in \cite[Section~3.5.1, Claim~2]{CGLSS2026}.  Use the new metric $d$
and witness $W$, and put $t=1-\delta$.  If $j$ was active and becomes
witness-inactive while $\ell$ remains witness-active, the
former active bid $\tau_j$, old phase value $z$, and new clipped bid
$\tau'_\ell$ satisfy
\[
 \begin{aligned}
 \tau_j&\ge z\ge t\,d(j,W)\\
 &\ge t\,d(\ell,W)-t\,d(j,\ell)\\
 &\ge\tau'_\ell-d(j,i_0)-d(\ell,i_0)
       \qquad(i_0\in F).
 \end{aligned}
\]
Subtracting $d(i_0,j)$ and taking positive parts shows that the
former active term dominates the new inactive dual term.  The scaled
distance $t\,d(j,W)$, rather than $d(j,W)$, is what the transfer
hypothesis supplies.  Together with clipping of surviving active bids,
this is the dual comparison; the payment tolerance is not added to the
dual inequality.

Witnesses and bids can be retained explicitly through the merge.
Deleting an earlier entry preserves the superset relation for an
inherited witness; a parameter transfer clips its inherited bids at
the new caps; and a newly completed entry uses its actual prefix.
These are polynomial-size certificates.  Fix deterministic LP and
tie-breaking choices in the source's completion subroutine when
comparing two interpolation endpoints.

\paragraph{Paid counts at interpolation endpoints.}
At a zero-length-copy endpoint, a regular facility and its free copy
are geometrically coincident but have different paid counts.  In
particular, the corresponding continuations can differ by one regular
facility; the geometric endpoint identities in
\cite[Section~3.5.2, Case~2; Section~3.5.3, Claim~5 and Case~3]{CGLSS2026}
must not be used as equalities of paid counts. Test each endpoint explicitly for the
target count $h$.  If an endpoint has count $h$, return it; otherwise
a change by one cannot pass from strictly below $h$ to strictly above
$h$.  This gives the needed integer bracketing for the interpolation
while preserving exactly $h$ paid regular facilities in the normal
branch.  No claim of exactly $h$ distinct projected facilities is used.

\paragraph{Disabling copies and bounding the phase count.}
Replace the fixed disabling length $10M$ used for
free copies in the source's Section~3.5 by $20NM$. This ensures that
the endpoint of a copy-length search has no effect except the copy's
label. Indeed, the no-overbidding invariant in
Lemma~\ref{imp:surplus-history} depends only on valid sequences.
For any active client $j$ at any atomic boundary
it implies $[z-d(i,j)]^+\le\lambda$ for every regular $i$, and hence
\[
 z\le M+\lambda\le Z:=M+8NM.
\]
An inactive client's scaled connection distance is at most its frozen
budget, which is at most $KZ$. Since $v<1/6$ and $N\ge3$,
\[
 KZ\le\frac{37}{36}(M+8NM)<10NM<(1-\delta)20NM.
\]
A copy of length $20NM$ therefore cannot inactivate an active client,
bind any bid cap at most $Kz$, or reduce an inactive client's
connection distance. These observations apply also to witness sets:
their scaled inactive connection distances are at most $z$, and their active bid caps
are at most $Kz$. The disabling endpoint consequently leaves all
subsequent openability tests and updates unchanged.

All these parameters have polynomial bit length. Their binary searches
remain polynomial. The bound $z\le Z$ and geometric phase growth give
$O_v(\log N)$ possible active phase indices, uniformly over the
search range. Multiplying a common phase bound by the source's at
most three free copies per phase gives the asserted uniform $L$.

\paragraph{Precision and parameter order.}
The total payment error in the normal branch is
\begin{equation}
 e=hN\eta\le N^2 2^{-N}.
 \label{imp:precision}
\end{equation}
For every fixed $v>0$ there is a fixed $N_v$ such that
$N^2 2^{-N}\le v$ for all $N\ge N_v$.  Every positive integral
comparison cost $P$ is at least one, so $e\le vP$ in that range.
The binary representation of $\eta$ has $O(N)$ bits.
Increasing the constant in $L=O_v(\log N)$ can place every
$N<N_v$ in the branch $k\le2L$, where the logarithmic recovery
procedure is used instead.  Thus constants in the local and
strictness arguments can be chosen first, then $v$, then the common
surplus bound and its fixed-size enlargement.  No dependence of
those earlier constants on $L$ is introduced.
